\input amstex
\documentstyle{amsppt}
\magnification 1200
\pageno 19
\define\eps{\epsilon}
\define\okk{{\Cal O}_K}
\define\off{\overline{F}}
\define\ovf{\overline{F}}
\define\unb{\underline{B}}
\define\unt{\underline{T}}
\define\unx{\underline{x}}
\define\qqq{\Bbb Q}
\define\ccc{\Bbb C}
\define\zzz{\Bbb Z}
\define\rrr{\Bbb R}
\define\aaa{\Bbb A}
\define\nnn{\Bbb N}
\define\gga{\Bbb G}
\define\oee{{\Cal O}_E}
\define\VV{\Cal V}
\define\HH{\Cal H}
\hsize 6.5 true in
\vsize 8.7 true in
\vskip .3in
\subheading{\S 2 Galois action on the coordinates:
explicit computations for $G_2$}
\vskip .1in
Let $G$ denote a connected semisimple $F$-quasi-split
group with no outer automorphisms.

If $T$ is a Cartan of $G$ defined over $F$ then
$\Gamma=Gal(\overline{F}/F)$ preserves
$T$ and the set $\Sigma$ of roots of $G$ with respect to $T$,
$\sigma\in \Gamma$ sending $U_\rho$ to $U_{\sigma\rho}$, where
$U_\rho$ is a root group as in \S1. Generally
speaking, $\Gamma$ will not preserve the set $\Delta$ of positive
simple roots with respect to a Borel $B$ containing $T$, unless
$B$ is defined over $F$. The character lattice 
$X:=X^*(T)=X^*(T/\overline{F})$ is a $\zzz[\Gamma]$-module and the action
of $\Gamma$ on $X_{\rrr}:=X^*(T)\otimes \rrr$ preserves the
set of root hyperplanes. Consequently, $\Gamma$ acts on the
set of Weyl chambers in $X_{\rrr}$ and preserves the positive 
Weyl chamber if $B$ is defined over $F$.
The action of $\sigma\in\Gamma$ on the set
${\Cal W}$ of Weyl chambers in $X^*(T_G)\otimes \rrr$
will be denoted $\sigma_T$.
\par
As in section 1, we fix a Borel $\underline{B}$ of $G$ 
defined over ${F}$, a Cartan 
$\underline{T}$ defined over $F$ contained in $\underline{B}$,
and a positive Weyl chamber $W_0$ with respect to the linear order 
on $\Sigma_G^+$ which makes all the roots occuring in $\underline{B}$
positive. Let $T_G$ denote our fixed Cartan of $G$ defined over $F$ 
and let $B_G$ denote a Borel containing
$T_G$. Note that, for any $\sigma\in\Gamma$ both $B_G$ and
$B_G^\sigma$ are Borels containing $T_G$, since $T_G$ is
defined over $F$. 
The Galois action on $\underline{T}$ induces a homomorphism 
$$
\rho:\Gamma\rightarrow Aut(X,\Delta,X^{\vee},\Delta^{\vee}),
\tag 2.1a
$$
where $X=X^*(\underline{T})$ denotes the root lattice, 
$X^{\vee}=X_*(\underline{T})$ denotes the cocharacter lattice,
$\Delta=\Delta(\underline{T},\underline{B})$ the simple roots of
$\underline{T}$ with respect to $\underline{B}$,
and $\Delta^{\vee}=\Delta^\vee(\underline{T},\underline{B})$
denotes the simple coroots.
This induces (\cite{Sp2}) a homomorphism
$$
\tilde{\rho}:\Gamma\rightarrow 
Aut(G,B,T,\{X_{\alpha}\}_{\alpha\in\Delta}\}),
\tag 2.1b
$$
where $X_\alpha$ denotes a root vector associated to $\alpha$
in the Lie algebra of $G$.
\par
The Borels
$\{{B}_G(W))\ |\ W\in {\Cal W}\}$ are precisely the Borel
subgroups of $G$ which contain ${T}_G$ and, hence, are preserved by the
Galois action since ${T}_G$ is defined over $F$. Define
a map $\sigma_T:{\Cal W}\rightarrow {\Cal W}$ by
$$
\sigma {B}_G(W)={B}_G(\sigma_T W),
$$
for $\sigma\in\Gamma$. This defines a continuous homomorphism
$\Gamma\rightarrow \Omega$,
$\sigma\longmapsto \sigma_T$, which factors through $Gal(\overline{F}/E)$
where $E/F$ denotes the splitting field of $T_G$.
We may write $B_G=\underline{B}^g$,
for some $g\in G(\overline{F})$ such that
$g^{-1}\cdot \sigma g\in T_G(\overline{F})$ for all $\sigma\in \Gamma$.
This action $\sigma_T$ on the chambers $W\in {\Cal W}$ induces an action,
also denoted $\sigma_T$, on the positive roots $\Sigma_G^+$ and hence on
the walls $ker(\rho)$ of $W$ to the walls
$ker(\sigma_T(\rho))$ of
$\sigma_TW$.

\par
We define an action of
$\Gamma$ on the set of regular stars, denoted $S_0$, by
$$
\sigma (B(W)):={B}_G(W)^{\sigma h}
=\sigma{B}(\sigma_T^{-1}W),
\tag 2.2
$$
where $B(W)=B_G(W)^h$.
The set of $F$-rational points of $S_0(\overline{F})$ 
under this $\Gamma$-action is the set
of regular stars $(B(W))_{W\in {\Cal W}}$ such that
$B(W)=B_G(W)^h$, with $\sigma(h)h^{-1}\in T_G(\overline{F})$
for all $\sigma\in\Gamma$.
\par

For $\gamma\in T_G(F)$, let
$$
\phi_\gamma^{-1}(F):=\phi^{-1}(\gamma)^\Gamma=
\{x\in X^0(F)\ |\ \phi(x)=\gamma\},
$$
where $\phi$ is defined following (1.2) above.
In fact, the projection map $\pi$ of \S 1.1
yields a set-theoretic bijection
$$
\phi_\gamma^{-1}(F)
\rightarrow \{g^{-1}\gamma g\ |\ g\in (T_G\setminus G)(F)\}
\cong (T_G\setminus G)(F),
\tag 2.3
$$
onto the stable conjugacy class of $\gamma$.


\par
The action of $\sigma\in\Gamma$ on a point 
$((B_-^{n_Wn})_{W\in{\Cal W}})\in
S(B_+)$, for some
$n,n_W\in N_+$, which we identify with the point
$((B(W))_{W\in{\Cal W}},n)\in S(B_+,B_-)\times N_+$,
is given by
$$
\sigma: S(B_+,B_-)\times N_+ \rightarrow S(B_+,B_-)\times N_+,
$$
$$
((B(W))_{W\in{\Cal W}},n)\longmapsto
((B'(W))_{W\in{\Cal W}},n'),
\tag 2.4a
$$
where $\sigma B_-=B_-^{c_\sigma}$,
$c_\sigma\in N_{Aut\ G}(T_+)(\off)=N_G(T_+)(\off)$, and
$$
B'(W)^{n'}=\sigma(B(\sigma_T^{-1}(W))^n)=
(\sigma B_-)^{\sigma(n_{\sigma_T^{-1}(W)}n)},
\tag 2.4b
$$
and
$$
n'=c_\sigma\sigma(n_{\sigma_T^{-1}(W_0)}n),\ \ B'(W_0)= B_- \ .
\tag 2.4c
$$
\par
How does the action of $\Gamma$ on $Y_1(B_+)$ affect the
point determined by $(b,(B_-^{n_W})_{W\in{\Cal W}})^{n}$?
Here $b\in B_-$, $n_W,n\in N_+(\ovf)$. 
Let $\alpha(W,\rho)\in\Delta$ denote the simple root associated to
a pair $(W,\rho)$ as above.
To be consistent with (2.4a), define an action of $\sigma\in\Gamma$
on $\Delta$, denoted by $\sigma_T^*$, as follows:
$$
(\sigma_T^*)^{-1}(\alpha(W,\rho))
=\alpha(\sigma_T^{-1}(W),\sigma_T^{-1}(\rho)).
\tag 2.5
$$
The action of $\Gamma$ on 
$S_1(B_+)\cong (S'(B_+,B_-)\times_{T_+}\aaa^\Delta)\times N_+$,
is given by
$$
\sigma:((B(W))_{W\in{\Cal W}},(z(\alpha))_{\alpha\in \Delta},n)\longmapsto
((B'(W))_{W\in{\Cal W}},(z'(\alpha))_{\alpha\in \Delta},n'),
$$
where $B'(W)$ and $n'$ are as above, and
$$
z'(\alpha)=\sigma(z((\sigma_T^*)^{-1}(\alpha))).
$$
By the formulas
for the Galois action on $S(B_+)$, we have
$$
\sigma:(b,(B_-^{n_W})_{W\in{\Cal W}})^{n}\longmapsto
(\sigma b,(\sigma B_-^{\sigma (n_{\sigma^{-1}_T (W)})}
)_{W\in{\Cal W}})^{\sigma(n_{\sigma^{-1}_T (W_0)}n)}.
\tag 2.6
$$
Here $b$ is as in (1.17) and
$$
n_W=\exp (z(W_r,\rho_r)X_{-\alpha_r})
\exp (z(W_{r-1},\rho_{r-1})X_{-\alpha_{r-1}})\cdots 
\exp (z(W_1,\rho_1)X_{-\alpha_1}),
$$
where $W_0,...,W_r=W$ is a gallery as in (1.9-10).
How does the action of $\Gamma$ on $S(B_+)$ affect the coordinates
$z(W,\rho)$? A point of $S(B_+,B_-)$ is determined by
$(z(W,\rho))_{(W,\rho)\ a \ pair}$, so by (2.2) and (2.6),
$$
\sigma:(z(W,\rho))_{(W,\rho)\ a\ pair}\longmapsto 
(z'(W,\rho))_{(W,\rho)\ a\ pair},
$$
where 
$$
z'(W,\rho)=\sigma(z(\sigma_T^{-1}(W),\sigma_T^{-1}(\rho)).
\tag 2.7
$$
Suppose now that we are on a patch $S(B_+,B_-)$ with
$T_+=B_+\cap B_-$ defined over $F$.
There is an $c_\sigma\in N_{Aut\ G}(T_+)(\off)=
N_G(T_+)(\off)$ such that
$\sigma B_-=c_\sigma B_-c_\sigma^{-1}$. We may then write
$$
\sigma:(b,(B_-^{n_W})_{W\in{\Cal W}})^{n}\longmapsto
(b',(B_-^{n'_W})_{W\in{\Cal W}})^{n'},
\tag 2.8
$$
where
$$
b'= \sigma({n_{\sigma^{-1}_T (W_0)}})c_\sigma b
c_\sigma\sigma({n^{-1}_{\sigma^{-1}_T (W_0)}})^{-1},
$$
$$
n'_W=c_\sigma
\sigma(n_{\sigma^{-1}_T (W)}n^{-1}_{\sigma^{-1}_T (W_0)})c_\sigma^{-1},
$$
$$
n'= c_\sigma \sigma(n_{\sigma^{-1}_T (W_0)}n).
$$
The action of $c_\sigma$ by conjugation commutes with the action of
$\sigma_T$ by permutation of the Borels in a star.
If $B_-$ is defined over $F$ then the
$W_0$-component of $(B_-^{n'_W})_{W\in{\Cal W}}$ is
$\sigma B_-=B_-$. In particular, if the point in $Y_1(B_+)$
determined by $(b,(B_-^{n_W})_{W\in{\Cal W}},
(z(\alpha))_{\alpha\in\Delta},{n_+})$ has the
coordinates $\{x(\rho)\ |\ \rho >0\}$ then the image of this point
under $\sigma$ has the coordinates $\{x'(\rho)\ |\ \rho >0\}$,
where $x'(\rho)$ is determined by
$$
\sigma({n_{\sigma^{-1}_T (W_0)}})c_\sigma b
c_\sigma\sigma({n^{-1}_{\sigma^{-1}_T (W_0)}})^{-1}
=b'=n_0t_0\prod_{\rho\in\Sigma^+}\exp (x'(\rho)X_\rho),
\tag 2.9
$$
for some $n_0\in N_+=<U_\rho\ |\ \rho>0>$,
where $\Sigma^+$ denotes the set of roots of $T_+$
which are positive with respect to $B_-$.
Using Chevalley's commutator formula, this formula for
$x'(\rho)$ will be used for explicit calculations below.
\proclaim{Lemma 2.10 \cite{H2, pf of Lemma 5.3}}
(a) Let $\alpha$ denote an 
arbitrary simple positive root of $G$, an arbitrary quasi-split group. If
$(u,(B(W))_{W\in {\Cal W}})\in E_{\Delta-\alpha}$ 
then every such Borel subgroup $B(W)$ lies in the same line $\ell_\beta$
in the Dynkin curve for $G$. Here
$E_{\Delta-\alpha}$ is a $1$-regular Richardson branch
of $\phi^{-1}(1)$ on $Y_C$ as in \S 3 below.

(b) The action of Galois on the points of $E_{\Delta-\alpha}$ 
is determined by the action on the lines in the Dynkin curve.
\endproclaim
\par
Now we restrict to the case $G=G_2$.

Recall the Dynkin curve for $G_2$ \cite{St}:
$$
$$
$$
$$
$$
$$
$$
$$
$$
$$
\vskip .4in
\noindent
Possible Galois actions on this diagram include:
\newline
(i) exchanging $\ell_\beta'$ with $\ell_\beta'''$, leaving 
$\ell_\beta''$ fixed, and flipping $\ell_\alpha$;
\newline
(ii) flipping $\ell_\beta'$, $\ell_\beta'''$, $\ell_\beta''$
and fixing $\ell_\alpha$;
\newline
(iii) cyclically permuting the lines $\ell_\beta'$, 
$\ell_\beta'''$, $\ell_\beta''$
and ``rotating'' $\ell_\alpha$.
\newline

Label and order the positive roots of the maximal torus 
$T_+\subset B_-$ of $G_2$ as
$$
\alpha \text{  (short) },\beta \text{  (long) },
\gamma=\alpha +\beta, \delta =2\alpha +\beta,
\epsilon =3\alpha +\beta,  \zeta=3\alpha +2\beta .
$$
We have
$$
\theta_\delta(r)\theta_\gamma(s)=\theta_\gamma(s)\theta_\delta(r)\theta_{\zeta}(3rs),
$$
so
$$
\split
\theta_{\eps}(x)^{\theta_{-\alpha}(z)}
&=\theta_{\eps}(x)\theta_\delta(c_1xz)\theta_\gamma(c_2xz^2)
\theta_\beta(c_3xz^3)\\
&=\theta_{\eps}(x)\theta_\beta(c_3xz^3)\theta_\delta(c_1xz)\theta_\gamma(c_2xz^2)\\
&=\theta_\beta(c_3xz^3)\theta_\gamma(c_2xz^2)\theta_\delta(c_1xz)\theta_{\eps}(x)
\theta_\zeta(c_4x^2z^3),
\endsplit
$$
where the $c_i$'s are computable constants. Likewise,
$$
\split
\theta_{\delta}(x)^{\theta_{-\alpha}(z)}
&=\theta_{\delta}(x)\theta_\gamma(d_1xz)\theta_\beta(d_2xz^2)\\
&=\theta_\beta(d_2xz^2)\theta_\gamma(d_1xz)\theta_{\delta}(x)\theta_\zeta(d_3x^2z),
\endsplit
$$
$$
\theta_{\gamma}(x)^{\theta_{-\alpha}(z)}
=\theta_{\gamma}(x)\theta_\beta(exz)=\theta_\beta(exz)\theta_{\gamma}(x),
$$
$$
\theta_{\beta}(x)^{\theta_{-\alpha}(z)}=\theta_{\beta}(x),
$$
and $\theta_{\alpha}(x)^{\theta_{-\alpha}(z)}=
\theta_{-\alpha}(y_\alpha)t\theta_{\alpha}(x)$, for some $y_\alpha$ depending on $x$, 
$z$, and some $t$ belonging to a
maximal torus.
\par
In general, if $\sigma\in\Gamma$ acts on 
$T_G(\overline{F})$ pointwise as as simple
reflection
$s_\rho$, some $\rho\in\Delta$, then we write $\sigma=\sigma_\rho$.
This amounts to taking $\sigma(n_{\sigma_T^{-1}(W_0)})=\theta_\rho (z)$, 
where $z=\sigma(z(\rho)z_1(W_0,\rho))$ \cite{H2, \S 3}. We will
compute the Galois action on the coordinate of a point $P$ in
$Y''(B_+,B_-)$, where we assume that $B_+\cap B_-=T_+$ is defined over
$F$ and $c_\sigma =1$ in (2.9). (We will allow $c_\sigma$ to
be a reflection below, see (2.16-17).)
\par
Taking $\sigma=\sigma_\alpha$, we have
$$
\prod_{\rho >0}\theta_\rho(x(\rho))^{\theta_{-\alpha}(z)}
=\theta_{-\alpha}(y_\alpha)
t\prod_{\rho >0}\theta_\rho(\phi_\rho^\alpha(\unx,z)),
\tag 2.11
$$
where $\unx =(x(\alpha),\cdots, x(\zeta))$,
$$
\split
\phi_\alpha^\alpha(\unx,z)
&=x(\alpha),\\
\phi_\beta^\alpha(\unx,z)
&=x(\beta)+ezx(\gamma)+d_2z^2x(\delta)+c_3z^3x(\eps),\\
\phi_\gamma^\alpha(\unx,z)
&=x(\gamma)+d_1zx(\delta)+c_2z^2x(\eps),\\
\phi_\delta^\alpha(\unx,z)
&=x(\delta)+c_2zx(\eps),\\
\phi_{\eps}^\alpha(\unx,z)
&=x(\eps),\\
\phi_\zeta^\alpha(\unx,z)
&=x(\zeta)+d_3zx(\delta)^2+c_4z^3x(\eps)^2+3c_2z^2x(\delta)x(\eps).
\endsplit
$$
Therefore, if $c_{\sigma_\alpha}=1$, we have
$$
\sigma_\alpha :x(\rho)\longmapsto \phi_\rho^\alpha(\unx,z),
\tag 2.12
$$
where $z=\sigma_\alpha(z(\alpha)z_1(W_0,\alpha))$.
\par
Take $\sigma =\sigma_\beta$ and note
$$
\theta_\gamma(r)\theta_\alpha (s)=\theta_\alpha (s)\theta_\gamma(r)
\theta_\delta(2rs)\theta_\epsilon(3rs^2)\theta_\zeta(3r^2s),
$$
$$
\theta_\beta(r)\theta_\alpha (s)=\theta_\alpha (s)\theta_\beta(r)
\theta_\gamma(rs)
\theta_\delta(r^2s)\theta_\epsilon(r^3s)\theta_\zeta(r^3s^2).
$$
We have, by Chevalley's commutator formula,
$$
\theta_\alpha(x)^{\theta_{-\beta}(z)}=\theta_\alpha(x),
$$
$$
\theta_\beta(x)^{\theta_{-\beta}(z)}=
\theta_{-\beta}(y_\beta)
t\theta_\beta(x),\ \ \ (t\in T_G ,\ 
\alpha(t)=1+xz,\  y_\beta:={-x^2z\over 1+xz}),
$$
$$
\split
\theta_\gamma(x)^{\theta_{-\beta}(z)}
&=\theta_\gamma(x)\theta_\alpha(czx)\\
&=\theta_\alpha(czx)\theta_\gamma(x)\theta_\delta(2czx^2)
\theta_\epsilon(3c^2z^2x^3)\theta_\zeta(3czx^3),
\endsplit
$$
$$
\theta_\delta(x)^{\theta_{-\beta}(z)}=\theta_\delta(x),
$$
$$
\theta_\epsilon(x)^{\theta_{-\beta}(z)}=\theta_\epsilon(x),
$$
$$
\theta_\zeta(x)^{\theta_{-\beta}(z)}=\theta_\zeta(x)\theta_\epsilon(dzx)
=\theta_\epsilon(dzx)\theta_\zeta(x).
$$
Therefore, 
$$
\prod_{\rho >0}\theta_\rho(x(\rho))^{\theta_{-\beta}(z)}
=t\prod_{\rho >0}\theta_\rho(\phi_\rho^\beta(\unx,z)),
\tag 2.13
$$
where
$$
\split
\phi_\alpha^\beta(\unx,z)
&=x(\alpha)(1+x(\beta)z)+cx(\gamma)z,\\
\phi_\beta^\beta(\unx,z)
&=x(\beta),\\
\phi_\gamma^\beta(\unx,z)
&=x(\gamma)+czx(\beta)x(\gamma),\\
\phi_\delta^\beta(\unx,z)
&=x(\delta)+2czx(\gamma)^2+c^2x(\beta)x(\gamma)^2z^2,\\
\phi_{\eps}^\beta(\unx,z)
&=x(\eps)+dzx(\zeta)+3c^2x(\gamma)^3z^2+c^3x(\beta)x(\gamma)^3z^3,\\
\phi_\zeta^\beta(\unx,z)
&=x(\zeta)+3zx(\gamma)x(\delta)+9czx(\gamma)^3+c^3z^3x(\beta)^2x(\gamma)^3.
\endsplit
$$
Therefore, if $c_{\sigma_\beta}=1$, we have
$$
\sigma_\beta :x(\rho)\longmapsto \phi_\rho^\beta(\unx,z),
\tag 2.14
$$
where $z=\sigma_\beta(z(\beta)z_1(W_0,\beta))$.
\par
Formulas (2.12) and (2.14) help to determine the set of $F$-rational
points of $Y_C$. We also need to compute the effect of conjugating $b$ by
$c_\sigma$. For this we recall a basic lemma \cite{Car}:

\proclaim{Lemma 2.15}
Let $\phi_\rho:SL(2)\rightarrow G$ be the
homomorphism associated to the root
$\rho$ and let $\theta_\rho:\gga_a\rightarrow U_\rho$ denote the root
group embedding, so 
$\theta_\rho(x)=\phi_\rho(\pmatrix 1&x\\ 0&1\endpmatrix)$. Let
$s_\rho:=s_\rho(1)$ where 
$$
s_\rho(x)=\theta_\rho(x)\theta_{-\rho}(-x^{-1})\theta_\rho(x)
=\phi_\rho(\pmatrix 0&x\\ -1/x&0\endpmatrix).
$$
Then 
$$
s_\rho \theta_\sigma (x)s_\rho^{-1}
=\theta_{w_\rho(\sigma)}(\eta_{\rho,\sigma}x),
$$
where $\eta_{\rho,\sigma}$ is as in \cite{Car, pp 94-96} and $w_\rho$ is 
the reflection of the Weyl group associated to $\rho$.
\endproclaim 

As a consequence of this, we may write
$$
\split
&s_\alpha\cdot \prod_{\rho>0}\theta_{\rho}(x(\rho))\cdot s_\alpha^{-1}\\
&=\theta_{-\alpha}(x'(\alpha))
\prod_{\rho>0, \rho\not=\alpha}\theta_{\rho}(x'(\rho)),
\endsplit
\tag 2.16
$$
where $x'(\alpha)=\eta_{\alpha,\alpha}x(\alpha)$,
$x'(\beta)=\eta_{\alpha,\epsilon}x(\epsilon)$,
$x'(\gamma)=\eta_{\alpha,\delta}x(\delta)$,
$x'(\delta)=\eta_{\alpha,\gamma}x(\gamma)$,
$x'(\epsilon)=\eta_{\alpha,\beta}x(\beta)$,
and
$$
x'(\zeta)=\eta_{\alpha,\zeta}x(\zeta)+
f_1x(\delta)x(\gamma)+f_2x(\beta)x(\epsilon),
$$
where $f_1,f_2$ are constants. Likewise, we may write
$$
\split
&s_\beta\cdot \prod_{\rho>0}\theta_{\rho}(x(\rho))\cdot s_\beta^{-1}\\
&=\theta_{-\beta}(x'(\beta))
\prod_{\rho>0, \rho\not=\beta}\theta_{\rho}(x'(\rho)),
\endsplit
\tag 2.17
$$
where 
$$
x'(\alpha)=\eta_{\alpha,\gamma}x(\gamma)+f_3x(\alpha)x(\beta)
$$
$x'(\beta)=\eta_{\alpha,\beta}x(\beta)$,
$x'(\gamma)=\eta_{\alpha,\alpha}x(\alpha)$,
$$
x'(\delta)=\eta_{\alpha,\delta}x(\delta)+f_4x'(\alpha)x'(\gamma),
$$
$x'(\epsilon)=\eta_{\alpha,\zeta}x(\zeta)$,
and
$$
x'(\zeta)=\eta_{\alpha,\epsilon}x(\epsilon).
$$
\enddocument
How do these coordinates relate to the coordinates on a unipotent orbit
of $G$? In general this does not appear to be an easy question, but if the
orbit is a Richardson orbit then the answer is quite simple. The discussion
below is motivated by \cite{H1, Remark V.2.2}.
\par
Let ${\Cal O}$ be the Richardson orbit associated to a parabolic $P$
as in \cite{St, Theorem 2,\S 3.9}. Write $P=MN_P$, so that $N_P^G$
contains ${\Cal O}$ as a dense subset, ${\Cal O}\cap N_P$ is dense in
$N_P$ and ${\Cal O}\cap N_P$ forms a single class under $P$. We have
dim${\Cal O}$ = dim$G$ - dim$M$. Let $u\in {\Cal O}\cap N_P$ and recall
$N_P^- P$ is dense in $G$ by the Bruhat decomposition. The $P$-orbit
$\{u^p\ |\ p\in P\}$ is equal to ${\Cal O}\cap N_P$ and is dense in $N_P$.
Therefore, the coordinates of $N_P$ yield coordinates for the $P$-orbit
of $u\in {\Cal O}\cap N_P$. The $G$-orbit of this $u$, that is
${\Cal O}$ itself, contains the $N_P^- P$-orbit of $u$ as an open 
subvariety. We may regard the $N_P^- P$-orbit of $u$ as the
$N_P^-$-orbit of ${\Cal O}\cap N_P$. By Corollary 1.6, for each
$\overline{n}\in N_P^-$, $B^{\overline{n}}$ is opposite to
$B^-$ and $B^{\overline{n}}\not= B$ unless $\overline{n}=1$.
Therefore, for our $u$ as above, $u^{\overline{n}}\notin {\Cal O}\cap N_P$.
From this it follows that the $N_P^-$-orbit of ${\Cal O}\cap N_P$
is isomorphic to $N_P^-$.
\proclaim{Lemma 1.21} If ${\Cal O}$ is the Richardson orbit associated to a 
parabolic $P$ then the coordinates on $N_P^- \times N_P$ serve as local
coordinates on ${\Cal O}$.
\endproclaim








